\chapter{Traumatologische Notfälle ---
\gAasRsN{RS~VII}}

\author{Josef Emhofer, Michael Motal, Michael Withofner}

\minitoc

\paragraph{Trauma}

\gDefinition{Als Trauma  bezeichnet man eine {\textperiodcentered}~Sch\"adigung,
{\textperiodcentered}~Verletzung, oder {\textperiodcentered}~Wunde, die
durch eine Gewalteinwirkung physikalischer oder chemischer Art
verursacht wurde.}

%Die Traumata lassen sich grob ind folgende Klassen Einteilen:
% 
% \begin{itemize}
%     \item Sch\"adel-Hirn-Trauma (SHT)
%     \item Thorax trauma
%     \item Abdominal trauma / Bauchtrauma
%     \item Beckentrauma
%     \item Verletzungen der Wirbels\"aule
%     \item Polytrauma
%     \item Verletzungen der Extremit\"aten
%     \item Verbrennungen
% \end{itemize}



%%%%%%%%%%%%%%%%%%%%%%%%%%%%%%%%%%%%%%%%%%%%%%%%%%%%%%%%%%%
%%%%%%%%%%%%%%%%%%%%%%%%%%%%%%%%%%%%%%%%%%%%%%%%%%%%%%%%%%%
%%%%%%%%%%%%%%%%%%%%%%%%%%%%%%%%%%%%%%%%%%%%%%%%%%%%%%%%%%%

\section{Einsatztaktische Überlegungen bei verunfallten Patienten}
Der Unfallmechanismus gibt wertvolle Hinweise auf wahrscheinliche bzw. mögliche Verletzungen des Patienten. Der Rettungssanitäter kann durch eine systematische Vorgehensweise bei der Patientenversorgung, die eine genaue Betrachtung des Verletzungsmechanismus beinhaltet, mithelfen, die Schäden eines Traumas begrenzt zu halten. Die wichtigsten Fragen dabei sind immer:
\begin{itemize}
\item Was ist passiert?
\item Wie ist der Patient verletzt worden?
\end{itemize}
Bei \gEs{Verkehrsunfällen} müssen bei der Einschätzung des Verletzungsbildes drei Faktoren beachtet werden:
\begin{itemize}
\item Deformierungsgrad des Fahrzeugs (Indiz für die involvierten Kräfte)
\item Deformierung von Teilen der Fahrzeugkabine (Hinweis für Aufschlagpunkt des Körpers im Fahrzeug)
\item Deformierung (Verletzungsmuster) des Patienten (Anzeichen dafür, welche Körperteile direkt aufgeprallt sind)
\end{itemize}
Für Unfälle mit Fahrzeugen gilt allgemein:
\begin{itemize}
\item Die Geschwindigkeit des Fahrzeugs vor dem Unfall spielt eine wichtige Rolle. Je schneller das Fahrzeug unterwegs war, desto schwerer kann die Verletzung sein. Dabei steigt die Bewegungsenergie des Fahrzeugs mit dem Quadrat der Geschwindigkeit an. Eine Verdoppelung der Geschwindigkeit hat daher eine Vervierfachung der Energie zur Folge, welche im Falle eines Unfalls frei wird und die Verformung verursacht.
\item Die Gewalteinwirkung auf den Patienten entspricht in ihrer Richtung i.d.R. der Einwirkung auf das Fahrzeug.
\item Wird ein Insasse aus dem Fahrzeug geschleudert, steigert sich das Risiko schwerer Verletzungen (SHT, Wirbelsäule, Thorax, Abdomen etc.) um 300 \% (Böhmer et al. 2006). Auch bei den im Auto verbliebenen Insassen ist mit schweren Verletzungen zu rechnen.
\item Wenn ein Insasse getötet oder eingeklemmt wird, ist bei den weiteren Insassen das Risiko schwerer Verletzungen erhöht.
\item Knieverletzungen bei einem Autounfall können Hinweise auf Oberschenkel- und Beckenfrakturen sowie Hüftgelenksverrenkungen sein.
\end{itemize}
Bei Frontalzusammenstöße sollten folgende Merkmale Alarmzeichen für schwere Verletzungen (Thoraxtrauma mit Herzkontusion und Pneumothorax, HWS-Trauma, Leber- und Milzruptur, Beckenfraktur) sein:
\begin{itemize}
\item Knieanstoß am Armaturenbrett,
\item Eingedrückte, leicht nach außen gewölbte Windschutzscheibe (ev. Spinnennetzmuster) durch Anstoß mit dem Kopf,
\item sichtbare Veränderungen des Radstandes (z. B. Rad zur Seite geknickt, Achsenverschiebung um mehr als 30 cm am Unfallfahrzeug),
\item Fahrzeugverformung um mehr als 50 cm.
\end{itemize}
Beim Seitenaufprall sind typische schwere Verletzungen: HWS-Verrenkung, Thoraxtrauma, Akzelerationstrauma der Aorta, seitenabhängig Leber- oder Milzruptur, Frakturen des Beckens. Bei einem Heckaufprall sollte man an HWS-Trauma (Schleudertrauma) denken und auch eine Frontalaufprallkomponente (Insasse wird nach vorne geworfen) mit einbeziehen (Lenkradkontusion).

Wird ein \gEs{Fußgänger von einem Fahrzeug (frontal) erfasst} ist generell mit schweren Verletzungen (vor allem Schädel-Hirn-Trauma, Wirbelsäule) zu rechnen. Zur besseren Einschätzung kann man folgende Faustregel (nach Böhmer et al. 2006) verwenden:
\begin{itemize}
\item bis 50 km/h: Aufschlagen des Kopfes auf die Kühlerhaube
\item bis 70 km/h: Aufschlagen des Kopfes auf die Windschutzscheibe
\item mehr als 70 km/h: der Fußgänger wird über das Dach geworfen
\end{itemize}
Lebensbedrohliche Verletzungen sind bereits möglich, wenn Radfahrer oder Fußgänger mit mehr als 30 km/h durch ein Fahrzeug erfasst werden. Erkennbar kann dies durch Bremsspuren, die bei trockener Fahrbahn länger als 10-20 m sind, sein.

Bei \gEs{Stürzen aus großen Höhen} (z. B. Leiter, Gerüst) ist stets an HWS-, Wirbelsäulen- und SHT zu denken. Bei Sturz auf die gestreckten Beine findet man typischerweise eine Fraktur der Fersenbeine und/oder eine Stauchungsfraktur der Wirbelsäule (Th12/L1), sowie Schienbeinplateau-Frakturen. Ab einer Fallhöhe von 6 m ist von einer Polytraumatisierung auszugehen.

Traumen, die durch spitze Gegenstände (Stich- und Pfählungsverletzungen) oder Schusswaffenmunition (Schussverletzungen) hervorgerufen werden, werden \gEs{penetrierende Verletzungen} genannt. Stichverletzungen durch Messer, Scheren oder ähnliche Gegenstände werden zumeist mit einer geringeren Geschwindigkeit zugefügt als Schusswaffenmunition und haben daher eine geringere sekundäre Verletzungsfolge als Schussverletzungen. Wie jeder andere Fremdkörper auch, sollte ein noch im Körper des Patienten befindliches Messer niemals präklinisch entfernt werden. Stichwunden können oft sehr klein, schwach blutend und daher harmlos aussehen. Um das Verletzungsausmaß richtig abschätzen zu können, sollten folgende Aspekte beachtet werden:
\begin{itemize}
\item Welche Körperteile sind betroffen?
\item Wie viele Einstichwunden gibt es?
\item Wie lang war die Klinge?
\item Wie war der Einstichwinkel? (Welche Organe könnten noch betroffen sein?)
\item Wurde nach dem Einstich das Messer (der Gegenstand) im Körper des Patienten umgedreht?
\end{itemize}

Aus allen einsatztaktischen Überlegungen ergibt sich ein Bündel an Erstmaßnahmen, welche Sie bei Unfällen durchführen müssen.
\begin{gMassnahmenEnv}
\begin{itemize}
\item Sichern Sie die Unfallstelle ab!
\item Verschaffen Sie sich einen Überblick (Wie viele Verletzte? ...) um ggf. weitere Einsatzkräfte nachzufordern!
\item Erheben Sie den Verletzungsmechanismus! Stellen Sie (sich) folgende Fragen:
	\begin{itemize}
	\item Aus welcher Richtung hat die Kraft auf den Patienten eingewirkt?
	\item Wie schnell war das Fahrzeug unterwegs? bzw. Aus welcher Höhe ist der Patient gefallen?
	\item War der Patient angegurtet?
	\item Welche besondere Verformungen sind am Fahrzeug sichtbar?
	\item Wurde der Patient aus dem Fahrzeug geschleudert?
	\end{itemize}
\end{itemize}
\end{gMassnahmenEnv}


\section{Wunden}\label{topic:}\index{}

\section{Wunden}

\subsection{Einteilung von Wunden}
Wunden lassen sich nach verschiedenen Gesichtspunkten einteilen. Eine Einteilung bezieht sich auf die physikalische Verletzungsursache:
\begin{itemize}
\item mechanische Wunden (mechanische Gewalteinwirkung, z.B. Schnitt, Stich, Schuss)
\item thermische Wunden (Einwirkung von Hitze oder Kälte: Verbrennungen, Erfrierungen)
\item chemische Wunden (Verätzungen)
\end{itemize}
Eine andere Einteilung gliedert die Wunden nach den vorrangig betroffenen Strukturen:
\begin{itemize}
\item Hautdefekt
\item Muskelverletzung
\item Gefäßverletzung
\item Organverletzung
\item Knochenverletzung
\end{itemize}
Die präklinisch am häufigsten verwendete Klassifizierung erfolgt nach dem Entstehungsmechanismus und der Wundbeschaffenheit:
\begin{itemize}
\item Schnittwunde (glatte Ränder)
\item Stichwunde (glatter Rand, meist klein aber tief)
\item Risswunde (Wundränder unregelmäßig und ev. zerfetzt)
\item Quetschwunde (durch Auftreffen stumpfer Gegenstände mit Aufplatzen der Haut; Platzwunde)
\item Rissquetschwunde (RQW; Kombination aus Riss- und Quetschwunde)
\item Schürf- und Kratzwunde
\item Bisswunde (Kombination aus Riss-, Stich- und Quetschwunde)
\item Ablederung (großflächiger Gewebsdefekt, z.B. Skalpierung, große Lappenwunden)
\item Amputation (Abtrennung eines Körperteils)
\item Pfählung/Einspießung (Eindringen und Verbleiben von spitzen aber z.T. unregelmäßigen Fremdkörpern)
\item Bluterguss (Hämatom)/Prellung (durch stumpfe Gewalteinwirkung)
\item chemische und thermische Wunden, durch Strahlung verursachte Wunden
\item Schusswunde
\end{itemize}

\gTempPicMissing{Bilder von Wunden, z.B. aus Gorgaß 2005 Seite 373}

\subsection{Wundversorgung}
Je nach stärke der Blutung versorgen Sie Wunden mit einem Wundverband, einem Druckverband oder in extremen Ausnahmesituationen durch eine Abbindung. Als Erstmaßnahme bevor ein Verband zur Verfügung steht können Sie mit einer sterilen Kompresse die Wunde abdrücken oder auch eine zuführende Arterie abdrücken. Die Techniken haben Sie im Erste Hilfe-Teil Ihrer Ausbildung bereits gelernt. An dieser Stelle sollen nur die wichtigsten \gEs {Grundregeln der Wundversorgung} wiederholt werden:
\begin{itemize}
\item Berühren Sie die Wunden nicht!
\item Arbeiten Sie immer mit Handschuhen und mit sterilem Material!
\item Wunden sollen grundsätzliche \gE{nicht} gereinigt werden! Sie dürfen nur lose Glassplitter oder ähnliches vorsichtig aus der Wunde entfernen.
\item Verwenden Sie \gE{keine} Desinfektionsmittel!
\item Belassen Sie Fremdkörper (Messer, Schere oder sonstige Pfählungsgegenstände) in der Wunde! Stabilisieren Sie diese Gegenstände mit geeignetem Fixationsmaterial (z.B. Mullbinden).
\item Versuchen Sie \gE{nicht} ausgetretene Strukturen (z.B. Gehirn, Eingeweide) zurück zu stopfen! Solche Organteile sollten mit steriler NaCl-Lösung feucht gehalten und anschließend steril mit einem Verbandstuch abgedeckt werden!
\end{itemize}

\subsection{Gefahren bei Wunden}
Jede Wunde hat (in unterschiedlichem Ausmaß) eine Auswirkung auf den Gesamtorganismus. Je nach Wundart, stärke der Gewalteinwirkung, betroffener Struktur und Sorgfalt bei der Wundversorgung gehen verschiedene Gefahren von Wunden aus:
\begin{itemize}
\item \gEs{Verletzung wichtiger Gewebestrukturen (Organe):} Je nach betroffener Körperregion treten unterschiedliche Verletzungsmuster auf. In den folgenden Abschnitten lernen Sie über spezielle traumatologische Notfälle wie z.B. Schädel-Hirn-Trauma, Throax-Trauma, Bauchtrauma etc.
\item \gEs{Starke Blutung:} Bei großflächigen Wunden oder Verletzung wichtiger Blutgefäße kann es zu starken Blutungen und infolge dessen zu einem Volumenmangelschock kommen. Blutstillung und Schockbekämpfung sind die wichtigsten Maßnahmen.
\item \gEs{Starke Schmerzen:} Treten im Zusammenhang mit einer Verletzung starke Schmerzen auf muss ein Notarzt für eine Schmerztherapie gerufen werden.
\item \gEs{Eintritt von Krankheitserregern (Infektionsgefahr):} Jede Wunde ist eine Eintrittsstelle für Krankheitserreger in unseren Organismus. Traumatologische Patient sollte daher über einen aufrechten Tetanus-Impfschutz verfügen (siehe Kapitel Hygiene)\footnote{Zur Verifizierung eines aufrechten Tetanus-Impfschutzes muss jede Wunde einem Arzt vorgestellt werden. Eine Impfung muss spätestens alle 10 Jahre aufgefrischt werden, bei entsprechenden Verletzungen schon früher.\cite{Böhmer6}}. Bei Bisswunden kann über den Speichel einer erkrankten Tieres Tollwut übertragen werden. Versuchen Sie bei Bisswunden daher immer zu erfragen, wer der Besitzer des Tieres ist und vermerken Sie dies auf Ihrem Einsatzprotokoll. Unter Umständen kann es notwendig sein die Polizei einzuschalten.
\end{itemize}

%\subsection{Mechanische
%Wunden}\label{topic:wunden-meschanische}\index{Wunden!mechanische}

\begin{figure}[H]
    \begin{center}
\includegraphics[width=5cm]{./images/stop.jpg}
\captionof{figure}{\gAuthNotesPicLegalNotOk{}{}Muster
\gSourcePicture{}}
\end{center}
\end{figure}




\begin{tabularx}{\linewidth}[H]{XXX}
\\\toprule
\multicolumn{3}{c}{{\gTabHeading{Diashow Muster}}}
\\\midrule
\includegraphics[width=\linewidth]{./images/stop.jpg}
\captionof{figure}{\gAuthNotesPicLegalNotOk{}{}Muster
\gSourcePicture{}} &
\includegraphics[width=\linewidth]{./images/stop.jpg}
\captionof{figure}{\gAuthNotesPicLegalNotOk{}{}Muster
\gSourcePicture{}} &
\includegraphics[width=\linewidth]{./images/stop.jpg}
\captionof{figure}{\gAuthNotesPicLegalNotOk{}{}Muster
\gSourcePicture{}} \\\midrule
\includegraphics[width=\linewidth]{./images/stop.jpg}
\captionof{figure}{\gAuthNotesPicLegalNotOk{}{}Muster
\gSourcePicture{}} &
\includegraphics[width=\linewidth]{./images/stop.jpg}
\captionof{figure}{\gAuthNotesPicLegalNotOk{}{}Muster
\gSourcePicture{}} &
\includegraphics[width=\linewidth]{./images/stop.jpg}
\captionof{figure}{\gAuthNotesPicLegalNotOk{}{}Muster
\gSourcePicture{}} \\\bottomrule
\end{tabularx}

%\subsubsection{Stichwunden}\label{topic:stichwunden}\index{Stichwunden}

\begin{gWideParEnv}
\begin{tabularx}{\linewidth}[H]{XXX}
\\\toprule
\multicolumn{3}{c}{{\gTabHeading{Diashow}}}
\\\midrule
\includegraphics[width=\linewidth]{./images/pneumothorax_stichverletzung_diskret.jpg}
\captionof{figure}{\gAuthNotesPicLegalOk{}{}Stichverletzng:
unauffällig. \gSourcePicture{Hauer}} &
\includegraphics[width=\linewidth]{./images/stichverletzung_oberarm.jpg}
\captionof{figure}{\gAuthNotesPicLegalOk{}{}Stichverletzung
am Oberarm, Schwellung. \gSourcePicture{Hauer}} &
{leer}
\\\bottomrule
\end{tabularx}
\end{gWideParEnv}


%\subsubsection{Schnittwunden}\label{topic:schnittwunden}\index{Schnittwunden}

\begin{figure}[H]
    \begin{center}
\includegraphics[width=5cm]{./images/pulsader_schnitt.jpg}
\captionof{figure}{\gAuthNotesPicLegalOk{}{}Schnittverletzung
Oberhalb der "`Pulsader"'.
\gSourcePicture{Hauer}}
\end{center}
\end{figure}

%\subsubsection{Rißwunden}\label{topic:risswunden}\index{Rißwunden}

%\subsubsection{Quetschwunden}\label{topic:quetschwunden}\index{Quetschwunden}

%\subsubsection{Rißquetschwunden
%(RQW)}\label{topic:rqw}\index{Rißquetschwunden}\index{RQW}

% \begin{gDiasEnv3}
% \includegraphics[width=\linewidth]{./images/wunde_rqw1.jpg}
% \captionof{figure}{\gAuthNotesPicLegalOk{}{}Rißquetschwunde
% vor der Versorgung im Spital.
% \gSourcePicture{Hauer}}
% &
% \includegraphics[width=\linewidth]{./images/wunde_rqw2.jpg}
% \captionof{figure}{\gAuthNotesPicLegalOk{}{}Rißquetschwunde
% nach der Versorgung im Spital.
% \gSourcePicture{Hauer}}
% &
% \end{gDiasEnv3}



\begin{figure}[H]
    \begin{center}
\includegraphics[width=5cm]{./images/wunde_rqw1.jpg}
\captionof{figure}{\gAuthNotesPicLegalOk{}{}Rißquetschwunde
vor der Versorgung im Spital.
\gSourcePicture{Hauer}}
\end{center}
\end{figure}

\begin{figure}[H]
    \begin{center}
\includegraphics[width=5cm]{./images/wunde_rqw2.jpg}
\captionof{figure}{\gAuthNotesPicLegalOk{}{}Rißquetschwunde
nach der Versorgung im Spital.
\gSourcePicture{Hauer}}
\end{center}
\end{figure}

%\subsubsection{Schußwunden}\label{topic:schusswunden}\index{Schußwunden}




%\subsection{Thermische Wunden}\label{topic:wunden-thermische}\index{Wunden!thermische}

%\subsubsection{Brandwunden}\label{topic:brandwunden}\index{Brandwunden}

%\subsubsection{Erfrierungen}\label{topic:erfrierungen}\index{Erfrierungen}

%\subsection{Chemische Wunden}\label{topic:wunden-chemische}\index{Wunden!chemische}

%\subsubsection{Verätzungen}\label{topic:veraetzungen}\index{Verätzungen}


%%%%%%%%%%%%%%%%%%%%%%%%%%%%%%%%%%%%%%%%%%%%%%%%%%%%%%%%%%%
%%%%%%%%%%%%%%%%%%%%%%%%%%%%%%%%%%%%%%%%%%%%%%%%%%%%%%%%%%%

%\subsection{Stichverletzungen}

%NIEMALS  HERAUSZIEHEN!!!!!!!!

%steril abdecken + fixieren

%verletzte Strukturen??, dient ev. als {\quotedblbase}St\"opsel{\textquotedblleft}??

%Fremdk\"orper soll nicht tiefer rutschen



%%%%%%%%%%%%%%%%%%%%%%%%%%%%%%%%%%%%%%%%%%%%%%%%%%%%%%%%%%%
%%%%%%%%%%%%%%%%%%%%%%%%%%%%%%%%%%%%%%%%%%%%%%%%%%%%%%%%%%%
%%%%%%%%%%%%%%%%%%%%%%%%%%%%%%%%%%%%%%%%%%%%%%%%%%%%%%%%%%%

\section{Sch\"adel-Hirn-Trauma (SHT)}
%Folge von Gewalteinwirkung auf den  Sch\"adel
\gDefinition{Gewalteinwirkung auf den Kopf, die zusätzlich zu den nicht immer vorhandenen Hautwunden und zu den Frakturen des knöchernen Schädels Funktionsstörungen und Verletzungen des Gehirns hervorrufen.\cite{Gorgass7}}
Die häufigsten Ursachen von Schädel-Hirn-Traumata (SHT) sind Verkehrsunfälle, Stürze und Gewaltverbrechen. Je nach Schwere und Mechanismus können verschiedene Formen des SHT auftreten.

\paragraph{Arten}

\begin{itemize}
\item Gedecktes SHT: ohne  Er\"offnung der harten Hirnhaut (Dura  mater)
\item Offenes SHT: mit  Er\"offnung der harten Hirnhaut (Dura  mater)
\end{itemize}

\paragraph{Begleitende Verletzungen}

\begin{itemize}
    \item Verletzung der Kopfschwarte  (Ri{\ss}quetschwunde)
    \item Frakturen von
    \begin{itemize}
        \item Sch\"adeldach
        \item Sch\"adelbasis (Blutung aus Ohren, Liquoraustritt  aus der Nase)
        \item Gesichtssch\"adel
        \item Unterkiefer
    \end{itemize}
\end{itemize}

\gCave{Die Schwere eines SHT lässt sich nur bedingt von den äußerlich sichtbaren Verletzungen ableiten! Oft werden im Krankenhaus Frakturen festgestellt, ohne dass äußere Wunden vorliegen. Eine sorgfältige Untersuchung ist daher bei jedem SHT Pflicht!}

\begin{gMassnahmenEnv}
    \paragraph{Spezielle Maßnahmen beim SHT --- Allgemein}
    
    \begin{itemize}
        \item Atmung: Freimachen der Atemwege (ggf. Absaugen), häufige Kontrollen der Atemwege und der Atmung
        \item Lagerung: leicht erhöhter Oberkörper, ca.
        30\,\textdegree , wenn keine Kontraindikation
        besteht.
        \item je nach Unfallmechanismus HWS-Immobilisation (Stifneck)
        \item Wundversorgung
        \gTxBeiVit $O_2$ 10-15 \gLMin
%        \begin{itemize}
%            \item Gedecktes SHT:
%            \begin{itemize}
%                \item Steriler Verband
%            \end{itemize}
%            \begin{itemize}
%                \item steriler Verband locker auf offene
%                Verletzungen
%                \item Knochen / Hirnteile NICHT 
%                zur\"uckstopfen /  entfernen!
%                \item Ausgetretene Hirnmasse feucht steril abdecken
%            \end{itemize}
%        \end{itemize}
    \end{itemize}
\end{gMassnahmenEnv}

\subsection{Gedecktes SHT}

Ein gedecktes Schädel-Hirn-Trauma ist ein SHT ohne Eröffnung der harten Hirnhaut. Je nach Schwere lassen wird das SHT in drei Grade eingeteilt.
 
\paragraph{Klassifikation}

\begin{itemize}
\item \gT{Gehirnersch\"utterung} -- SHT 1\textdegree --
Contusio cerebri
\item \gT{Gehirnprellung} -- 2\textdegree -- Contusio
cerebri
\item \gT{Gehirnquetschung} -- SHT 2\textdegree 
\end{itemize}




\subsubsection{Gehirnersch\"utterung (Commotio (cerebri) --
SHT 1\textdegree}
\index{Gehirnersch\"utterung}\index{Commotio (cerebri)}

\gMarried{Die Gehirnerschütterung ist die leichteste Form des
Schädelhirntraumas. Leit Symptom ist ein Trauma mit
anschließender kurzer Bewusstlosigkeit. Weiters kann es zu
flüchtigen neurologischen Ausfällen kommen, es kann auch eine
retrograde Amnesie (Gedächtnisverlust betreffenden Vorgänge
oder dem Unfall) bestehen. Eventuell lag der Patient über
Übelkeit erbrechen und Schwindel. Die Gehirnrschütterung
sieht keine bleibenden Schäden nach sich.}
{
\begin{itemize}
    \item  Bewu{\ss}tlosigkeit /Bewu{\ss}tseinstr\"ubung {\textless}  1h
    \item \"Ubelkeit, Erbrechen, Schwindel
    \item evt. fl\"uchtige neurologische Ausf\"alle
    \item retrograde  Amnesie  ({\quotedblbase}kann sich an den
    Unfallhergang nicht erinnern{\textquotedblleft} )
    \begin{itemize}
        \item im CT/MR [${\rightarrow}$R{}-\pageref{bkm:LexComputertomographie}]
        keine nachweisbaren Sch\"aden
        \item keine  bleibenden Sch\"aden
    \end{itemize}
\end{itemize}}

%\begin{gMassnahmenEnv}
%    \paragraph{Spezielle Maßnahmen beim SHT ---Gehirnerschütterung}
%    \begin{itemize}
%        \item 
%    \end{itemize}
%\end{gMassnahmenEnv}


\subsubsection{Gehirnprellung (Contusio (cerebri) --
SHT 2\textdegree}

\index{Gehirnprellung}\index{Contusio (cerebri)}

\gMarried{Die Gehirnprellung ist eine schwerere Form des Schädel-Hirntraumas. 
Leitsymptom ist die \gEs{Bewusstlosigkeit größer
 als 1 h}, beziehungsweise die Bewusstseinseintrübung >24 h.
 Mittes weiterführender Diagnostik im Spital lassen sich die
 \gEs{Hirnschädigungen} nachweisen. Es kam zur
 Nervenverletzungen kommen, dies kann bis zu vegetativen Störungen führen.
 Bleibende Schäden sind möglich.}
 {
 \begin{itemize}
    \item Bewu{\ss}tlosigkeit  {\textgreater}  1h
    \item Bewu{\ss}tseinstr\"ubung {\textgreater}24h
    \item nachweisbare Gehirnsch\"adigung / bleibende  Ausf\"alle
    \item ev. Gehirnnervenabrisse 
    \item  vegetative St\"orungen
\end{itemize}}

%\begin{gMassnahmenEnv}
%    \paragraph{Spezielle Maßnahmen beim SHT --- Gehirnprellung}
%    \begin{itemize}
%        \item 
%    \end{itemize}
%\end{gMassnahmenEnv}




\subsubsection{Gehirnquetschung -- SHT 3{\textdegree}}

\index{Gehirnquetschung}

\gMarried{Die Gehirnquetschung ist eine sehr schwere Form des
Schädelhirntraumas. Es kommt zu einer \gEs{Bewusstlosigkeit
länger als 24 h} und zu teils sehr ausgeprägten
Begleitsymptomen wie zum Beispiel Hirndruckzeichen, Atem und
Kreislaufstörungen bis hin zum Atemstillstand durch ein Klemm
des Atemzentrums. 

Es kann weiteres zu intrakranielle
Blutungen kommen (epidural, subdural, subarachnoidal).
Außerdem kommt es zu verschiedenen neurologischen Ausfällen
und Symptomen wie zum Beispiel in einen oder auch zerebrale
Kampfanfälle. }
{
\begin{itemize}
    \item Bewu{\ss}tlosigkeit  {\textgreater}  24h
    \item Atmungs-/ Kreislaufst\"orungen
    \item Hirndruck symptomatik
    \item evt. Einklemmung des Atemzentrums ${\rightarrow}$
    zentraler Atemstillstand
    \item Behinderung der Hirndurchblutung
    \item Hirnblutungen
%     \begin{itemize}
%         \item epidural
%         \item subdural
%         \item subarachnoidal
%     \end{itemize}
    \item L\"ahmungen
    \item Zerebrale Kr\"ampfe
\end{itemize}}

%\begin{gMassnahmenEnv}
%    \paragraph{Spezielle Maßnahmen beim SHT --- Gehirnquetschung}
%    \begin{itemize}
%        \item 
%    \end{itemize}
%\end{gMassnahmenEnv}

Mit \gEs{bleibenden Schäden} und Langzeitfolgen ist
zu rechnen.

\begin{gMassnahmenEnv}
    \paragraph{Spezielle Maßnahmen beim SHT --- gedecktes SHT}
    \begin{itemize}
        \item steriler Verband (steriles Abdecken)
        \item Knochen/Fremdkörper nicht entfernen/zurückstopfen
    \end{itemize}
\end{gMassnahmenEnv}

\gFazit{Das Leitsymtptom eines jeden SHT ist die
\gEs{Bewusstseinsstörung}.}


\subsection{Offenes SHT}

Hinweise auf ein offenes SHT sind:

\begin{itemize}
    \item gro{\ss}e RQW mit tast-/sichtbaren  Knochensplittern
    \item Hirnaustritt (schlechte Prognose)
    \item Blutungen aus Nase/Ohren
    \item Liquoraustritt, oft mit Blut vermischt
    ({\quotedblbase}Tupfertest{\textquotedblleft})
\end{itemize}

\begin{gMassnahmenEnv}
    \paragraph{Spezielle Maßnahmen beim SHT --- gedecktes SHT}
    \begin{itemize}
        \item ausgetretene Hirnmasse feucht steril abdecken und locker befestigen
    \end{itemize}
\end{gMassnahmenEnv}


\subsection{Intrakranielle Blutungen}

\paragraph{\gT{Luzides Intervall} (lichtes Intervall)}

Die Symptome können eventuell erst Stunden nach der
Verletzung einsetzen. Dazwischen ist der Patient unauffällig.
Diese Zeitspanne wird Luzides Intervall genannt und ist
sehr gefährlich da man den Patienten leicht falsch
einschätzen kann.

\gMarried{Intrakranielle Blutungen können sowohl in Folge eines SHT als auch spontan (z.B. durch Platzen eines Aneurysmas) auftreten. Aufgrund der inneren Blutung steigt der Druck im Gehirn an, im Extremfall kann es zu einer Hirnstammeinklemmung kommen. Patienten mit intrakranieller Blutung sind vital bedroht! Beim SHT ist zu beachten, dass die Symptome oft auch erst Stunden nach dem Unfall auftreten können - siehe luzides Intervall.}
{\begin{itemize}
    \item stärkste Kopfschmerzen (langsam oder plötzlich einsetzend)
    \item Bewusstseinsstörungen bis Bewusstlosigkeit
    \item Pupillendifferenz, verlangsamte Pupillenreaktion
    \item ev. Lähmungen oder Krämpfe
    \item ev. Atemnot, unregelmäßige Atmung und Kreislaufstörung
    \item ev. Druckpuls (Bradykardie mit hohen Blutdruckwerten als Folge des Hirndrucks)
    \item Übelkeit, Erbrechen
\end{itemize}}

%\gAuthNotes{Zusammenfassen, keine Unterteilung in Sub-,
%Epi- und Aubarachnoidalblutung, nur mehr "`Hirnblutung"'}


%\subsubsection{Epidurale Blutung} \index{Epidurale Blutung}

%\gDefinition{Blutung zwischen
%der harten Hirnhaut und dem Schädelknochen.}

%\gMarried{Bei der epiduralen Blutung kommt es zu einer Blutung zwischen
%der harten Hirnhaut und dem Schädelknochen, bedingt durch
%den Riss einer Hirnhaut-Arterie. Durch die Blutung steigt
%der Hrindruck, es kommt dabei zu \gE{Hirndruckzeichen}.}
%{\begin{itemize}
%    \item Ri{\ss} einer Meningealarterie
%    \item Einblutung zwischen Kalotte und Dura mater
%    \item {\textquotedblleft}lichtes{\textquotedblright} (luzides )
%    Intervall: einige Stunden
%\end{itemize}}



%\subsubsection{Subdurale  Blutung} \index{Subdurale Blutung}

%\gDefinition{Blutung
%unterhalb der harten Hirnhaut und oberhalt der weichen
%Hirnhaut.}

%\gMarried{Bei der subduralen Blutung kommend es zu einer ein Blutung
%unterhalb der harten Hirnhaut. Der Verlauf ist oft
%schleichend und kann sich über mehrere Tage erstrecken. Die
%genaue Unterscheidung zwischen epiduraler und subduraler
%Blutung ist präklinisch nicht möglich.}
%{\begin{itemize}
%    \item schleichender (!) Verlauf: oft \"uber Tage
%    \item auch ohne  Trauma: Hypertoniker, Alkoholiker
%\end{itemize}}



%\subsubsection{Subarachnoidale Blutung -- SAB}
%\index{Subarachnoidale Blutung} (\index{SAB}SAB)

%\emph{Abkz.: \gT{SAB}}

% TODO test

%\gDefinition{Bei der subarachnoidalen Blutung kommt es zu einer
%Blutung unterhalb der Spinnengewebeshaut (Arachnoidea).} 

%\gMarried{Dafür gibt es
%unterschiedliche Ursachen, sie kann sowohl durch ein Trauma
%ausgelöst werden als auch spontan entstehen. Bei der
%\gE{spontanen} Subarachnoidalenblutung platzt typischerweise
%ein bereits bestehendes kleines Aneurysma eines Hirngefäßes und
%verursacht die Blutung. Dies kann aus heiterem Himmel
%passieren.

%\gFazit{Die subarachnoidalen Blutung kann durch ein Trauma
%oder spontan entstehen. Sie kann ein traumatologischer oder
%auch neurologischer Notfall sein.}

%Leitsymptom ist der plötzliche, schwere,
%\gEs{"donnerschlagartige" Kopfschmerz}, welcher mit
%Hirndruckzeichen und einem oft raschen Verfall des
%Patienten einhergeht. Der Patient ist akut \gE{vital
%bedroht!}}
%{
%\begin{itemize}
%    \item Trauma oder spontan
%    \item z.B. geplatztes Aneurysma, Br\"uckenvenenabri{\ss} bei
%    Decelerationstrauma
%    \item Klassisch: pl\"otzlicher, schwerer,
%    {\quotedblbase}donnerschlagartiger{\textquotedblleft} Kopfschmerz
%    \item evt. Herd symptomatik, schneller
%    Verfall,
%\end{itemize}}





\subsection{Spezielle Diagnostik beim SHT}

Grundsätzlich ist beim Schäde-Hirn-Trauma die übliche
Diagnostik druchzuführen, die besonders beim SHT wichtigen
Punkte sind hier jedoch nochmal beschrieben.

Beurteilung von:
\begin{itemize}
	\item Vitalfunktionen
	\item Neurologie (Neurocheck)
	\item Verletzungen (Traumacheck)
	\item Unfallmechanismus (Anamnese)
\end{itemize}

\paragraph{Neurologische Beurteilung}

Die neurologische Beurteilung ist beim SHT natürlich
besonders wichtig und besonders genau durchzuführen.
Auffälligkeiten können Hinweise auf eine direkte
Hirnschädigung oder eine Hirnblutung sein. 

\begin{description}
\item[Bewu{\ss}tseinslage] klar/getr\"ubt/bewu{\ss}tlos
\item[Motorik] 
\item[Pupillenweite] Hirndruckzeichen
\item[Pupillenreaktion] Hirndruckzeichen
\end{description}

\paragraph{Traumacheck} Ein Traumackeck ist obligat
und sollte keiner weiteren Erwähnung bedürfen.



\gAuthNotes{GABS: Die einzelnen Massnahmen in Bezug auf die
unterschiedlichen SHTs gehören definiert.}






\begin{gMassnahmenEnv}
    \paragraph{Spezielle Maßnahmen beim SHT --- Gehirnerschütterung}
    \begin{itemize}
        \item 
    \end{itemize}
\end{gMassnahmenEnv}





\begin{itemize}
\item (Gefahr: hohe Querschnittl\"ahmung)

\end{itemize}
\begin{itemize}
\item Sch\"adel/HWS/BWS bilden eine Linie

\end{itemize}
\begin{itemize}
\item Jedes SHT ist im Spital abzukl\"aren! (klinisch und  bildgebend)

\end{itemize}

\gCave{\gEs{Jedes SHT ist im Spital abzukl\"aren!} Es
besteht die Gefahr dass sich der Patient gegenw\"artig in der
\gEs{"`luciden Phase"'} 
befindet und sich sein Zustand erst in einigen Stunden, dann aber schlagartig verschlechtert. Im Zweifelsfall wird der
Patient dort zur  Beobachtung (24h lang) aufgenommen.}



%%%%%%%%%%%%%%%%%%%%%%%%%%%%%%%%%%%%%%%%%%%%%%%%%%%%%%%%%%%
%%%%%%%%%%%%%%%%%%%%%%%%%%%%%%%%%%%%%%%%%%%%%%%%%%%%%%%%%%%
%%%%%%%%%%%%%%%%%%%%%%%%%%%%%%%%%%%%%%%%%%%%%%%%%%%%%%%%%%%


\section{Wirbels\"aulentrauma}

\gAuthNotes{KOCH: Überarbeiten, evt etwas kürzen.}

\paragraph{Typische Unfallmechanismen}

\gDefinition{Gewalteinwirkung auf die Wirbelsäule, die zur Verschiebung oder zur Fraktur von Wirbeln mit oder ohne Rückenmarkschädigung führt.\cite{Gorgass7}}

\begin{itemize}
    \item Sturz aus gro{\ss}er H\"ohe (> 3 m)
    \item Motorradunfall
    \item eingeklemmte Personen bei VU
    \item Schlag auf Kopf (Stauchung der HWS)
    \item Schleuderbewegung bei VU, wenn Patient nicht angeschnallt war oder Kopfstütze fehlt
    \item Suizidversuch durch Erhängen
\end{itemize}

\paragraph{Gefahren}

\begin{itemize}
    \item Kompression des R\"uckenmarks (aufgrund von Einblutungen ins  R\"uckenmark und dadurch bedingtes Ödem) 
    \item Durchtrennung des R\"uckenmarks
    \item L\"ahmungen (Parese)
    \item Sensibilit\"atsst\"orungen (Par\"asthesie)
\end{itemize}

\begin{gMassnahmenEnv}
    \paragraph{Spezielle Maßnahmen beim Wirbelsäulentrauma --- Allgemein}
    
    \begin{itemize}
        \item bei ansprechbaren Patienten: immer HWS-Immobilisation (Stifneck)
        \item Rettung mit Schaufeltrage oder KED-System
        \item Lagerung auf/Schienung mit Vakuummatratze (wenn Schocklagerung notwendig, Tragentisch schräg stellen)
        \item ev. Sauerstoffgabe
        \item bei Querschnittslähmung: Kennzeichnung der Sensibilitätsgrenze mit Angabe der Uhrzeit
        \item ev. Wundversorgung
        \item ev. Schockbekämpfung
        \item regelmäßige Kontrolle der Atmung
        \item Erheben und Dokumentieren des Unfallmechanismus
    \end{itemize}
\end{gMassnahmenEnv}

\subsubsection{Schleudertrauma, Peitschenschlag}

Schleudertraumen kommen bei Auffahrunf\"allen vor, insbesondere dann, wenn der Patient nicht angeschnallt war oder die Kopfstütze fehlt. Es kommt zu einer Zerrung des Halsmarks, welche sich durch
\begin{itemize}
	\item Schmerzen im HWS-Bereich
	\item Kopfschmerzen
	\item Schwindel, Tinnitus
	\item ev. neurologische Ausfälle (in schweren Fällen)
\end{itemize}

\subsubsection{Querschnittssydrom}

\gMarried{Kommt es zu einer Druckerhöhung (z.B. aufgrund einer Blutung) im Rückenmark oder zu einer Durchtrennung des Rückenmarks können Symptome eines Querschnitts auftreten. Ein sorgsamer und vorsichtiger Umgang bei der Rettung, der Untersuchung und dem Transport des Patienten ist unbedingt notwendig.}{
\begin{itemize}
	\item Ausfall von Motorik und/oder Sensibilität
	\item Blutdruckabfall aufgrund fehlender Gefäßregulation
	\item schlaffe oder spastische Lähmung der Extremitäten
	\item Inkontinenz aufgrund von Blasen-/Darmlähmung
\end{itemize}
}

\subsubsection{Spinaler Schock}

\gMarried{Aufgrund des Unfalls kommt es zu einer plötzlichen Durchtrennung des Rückenmarks, wodurch die Gefäßregulation (und somit der Gefäßtonus) versagt. Somit sinkt der Blutdruck rasch ab, es kommt zu einer lebensbedrohlichen Kreislaufstörung.}{
\begin{itemize}
	\item deutlicher Blutdruckabfall aufgrund fehlender Gefäßregulation
	\item Querschnittslähmung	
	\item Atemlähmungen (bei hohem Querschnitt (im Bereich der HWS))
\end{itemize}
}

%%%%%%%%%%%%%%%%%%%%%%%%%%%%%%%%%%%%%%%%%%%%%%%%%%%%%%%%%%%
%%%%%%%%%%%%%%%%%%%%%%%%%%%%%%%%%%%%%%%%%%%%%%%%%%%%%%%%%%%
%%%%%%%%%%%%%%%%%%%%%%%%%%%%%%%%%%%%%%%%%%%%%%%%%%%%%%%%%%%

\section{Thoraxtrauma}

\gDefinition{Verletzung des Brustkorbes und/oder der darin
enthaltenen Organe.}

\paragraph{Thoraxtrauma -- allgemeine  Symptome}

\begin{itemize}
    \item Dyspnoe
    \begin{itemize}
        \item Tachypnoe
        \item Zyanose
    \end{itemize}
    \item atemabh\"angige  Schmerzen,  Inspirationsstop
    \item ev. Hautemphysem (Luftansammlung in der Haut, typisches
    {\quotedblbase}Knistern{\textquotedblleft})
    \item Prellmarken (z.B. Lenks\"aule von PKW)
    \begin{itemize}
        \item ={\textgreater} gr\"undlicher Traumacheck!
    \end{itemize}
\end{itemize}



\paragraph{Thoraxtrauma -- spezifische  Symptome}

\begin{itemize}
    \item instabiler Thorax
    \begin{itemize}
        \item Sternumfraktur
        \item Serienrippenfraktur (zus. paradoxe Atmung --  Brustkorb senkt beim
        Einatmen)
    \end{itemize}
    \item Dyspnoe + Prellmarken +  Bluthusten
    \begin{itemize}
        \item Lungenprellung: Gasaustausch [F0EA?] , da  Lungenabschnitte
        gesch\"adigt
    \end{itemize}
    \item Brustschmerz, Schock
    \begin{itemize}
        \item Rhythmusst\"orungen ${\rightarrow}$ Verletzung des  Herzens
        \item schwerer Schock, schneller Verfall ${\rightarrow}$ Riss der  Aorta
        \item Blut im Schwall aus dem Mund ${\rightarrow}$
        Speiser\"ohrenverletzung
    \end{itemize}
\end{itemize}

\subsubsection{Pneumothorax}\label{topic:pneumothorax}\index{Pneumothorax}

\begin{center}
 [Warning: Image ignored] % Unhandled or unsupported graphics:
%\includegraphics[width=7.086cm,height=10.686cm]{AASdev20090228-img91.png}
\gTempPicMissing{Pneumothorax img91}
\end{center}

\begin{itemize}
    \item Lungen- (innerer) oder  Brustwandverletzung (\"au{\ss}erer)
    \begin{itemize}
        \item Luft gelangt in den Pleuraspalt
        \item Unterdruck aufgehoben
        \item Lunge kollabiert auf betroffener Seite
    \end{itemize}
    \item "`unkomplizierter"' Pneumothorax
    \begin{itemize}
        \item Luft kann ungehindert in Pleuraspalt und  wieder heraus
        \item {\quotedblbase}nur{\textquotedblleft} Ausfall des betroffenen
        Lungenfl\"ugels
    \end{itemize}
    \item Ventil-, \index{Spannungspneumothorax}Spannungspneumothorax
    \begin{itemize}
        \item Luft kann zwar hinein, aber nicht wieder  heraus
        \item Druck auf betroffener Seite steigt ${\rightarrow}$  Organe (Lunge,
        Herz) werden auf die andere Seite gedr\"uckt
        \begin{itemize}
            \item gesunde Seite wird verdr\"angt + Abknickung  gro{\ss}er
            Gef\"a{\ss}e!
        \end{itemize}
    \end{itemize}
\end{itemize}

\begin{center}
 [Warning: Image ignored] % Unhandled or unsupported graphics:
%\includegraphics[width=5.811cm,height=4.343cm]{AASdev20090228-img92.png}
\gTempPicMissing{Spannungspneu img92}
\end{center}
Therapie

Pneumothorax, unkompliziert:

\begin{gMassnahmenEnv}
    \paragraph{Spezielle Maßnahmen --- Pneumothorax}
    \begin{itemize}
        \item \gTxMassVitBes
        \begin{itemize}
            \item $O_2$: 10\,-\,15\,l
            \item Lagerung: Oberkörper hoch. Wenn stabile
            Seitenlage notwendig ist dann Lagerung auf die
            \gE{verletzte} Seite, damit die unverletzte
            Lunge  nicht unnötig belastet wird.
        \end{itemize}
        \item Bei offener Thoraxverletzung: \gEs{Nicht
        luftdicht verbinden!} Es kann sonst ein
        Spannungspneumothorax entstehen. Steril abdecken. 
    \end{itemize}
\end{gMassnahmenEnv}

\gZ{Der Notarzt kann mittels Drainage einen Pneumothorax
entlasten.}

\gAuthNotes{KOCH: 

Es fehlt:

Serienrippenfraktur erklären
Paradoxe Atmung
Gefahr von Blutverlust
Therapie

Ev. auch kurz etwas über STernumfraktur schreiben?}

\section{Bauchtrauma}
\index{Bauchtrauma}

\begin{itemize}
    \item offen (perforierend),
    \begin{itemize}
        \item  Er\"offnung der  Bauchh\"ohle
        \item Stich-, Pf\"ahlungsverletzung
        \item Schnittverletzung, {\quotedblbase}Platzbauch{\textquotedblleft}
    \end{itemize}
    \item geschlossen (stumpf)
    \begin{itemize}
        \item Schlag, Tritt
        \item Prellung durch Lenkstange bei PKW-Unfall
    \end{itemize}
\end{itemize}




\subsection{Offenes Bauchtrauma}

\begin{center}
 [Warning: Image ignored] % Unhandled or unsupported graphics:
%\includegraphics[width=6.543cm,height=4.839cm]{AASdev20090228-img93.png}
\gTempPicMissing{Offenes Bauchtrauma img93}
\end{center}
\begin{itemize}
\item ev. Austritt von Darmschlingen

\end{itemize}

\begin{gMassnahmenEnv}
    
    \paragraph{Spezielle Maßnahmen beim offenen Bauchtrauma}
    \begin{itemize}
        \item \gTxMassVitBes
        \begin{itemize}
            \item Lagerung: Bauchdecke entspannen , Knierolle
        \end{itemize}
        \item Darmteile nicht  zur\"uckstopfen ${\rightarrow}$ sonst
        Darmri{\ss}/Darmverschlu{\ss}
        \item steril abdecken, mit Ringerl\"osung feuchthalten
        \begin{itemize}
            \item Gefahr des Eindringens von Keimen  (Peritonitis/Sepsis!)
        \end{itemize}
    \end{itemize}
\end{gMassnahmenEnv}


\subsection{Geschlossenes Bauchtrauma}

\paragraph{Gefahren}

\begin{itemize}
    \item Gef\"a{\ss}verletzungen
    \begin{itemize}
        \item Aorta, Vena cava,...
        \item massive innere Blutungen ${\rightarrow}$ Volumenmangelschock
    \end{itemize}
    \item Zerrei{\ss}ung  von Organen (\gEs{Milz},
    \gEs{Leber})
    \begin{itemize}
        \item massive \gEs{innere Blutungen} ${\rightarrow}$
        Volumenmangelschock
    \end{itemize}
    \item Zerrei{\ss}ung  von Hohlorganen (Magen, Darm,Galleng\"ange )
    \begin{itemize}
        \item Austritt von Speisebrei, Magensaft, Stuhl, Gallensekret
        \begin{itemize}
            \item Bauchfellentz\"undung (Peritonitis) nach Std. bis  Tagen
            \item Sepsis (hohe Lethalit\"at), septischer Schock
        \end{itemize}
    \end{itemize}
\end{itemize}

\paragraph{Symptome}

\begin{itemize}
    \item starke  Bauchschmerzen
    \item brettharte  Bauchdecke  (Abwehrspannung, D\'efense )
    \item Leitsymptom: Schockzeichen
    \item beim Bewu{\ss}tlosen schwierig zu  diagnostizieren
    \begin{itemize}
        \item gr\"undlich nach \index{Prellmarken}\gEs{Prellmarken}
        suchen!
    \end{itemize}
\end{itemize}

\gFazit{Immer nach \gEs{Prellmarken} suchen. \gEs{Patienten
entkleiden!}}


\subsubsection{Milzruptur}\index{Milzruptur}

Die Milz ist ein teuflisches Organ. Sie besitzt rund um
das gut durchblutete Gewebe eine Bindegewebs-Kapsel.

Das
Schlimme daran: Das Gewebe kann reißen, die Kapsel bleibt
jedoch intakt. Es kommt zuerst nur zu einer geringfügigen
Einblutung, dem Patienten geht es zunächst gut. Nach
einiger Zeit reißt dann die Kapsel und es kommt zu einer
ungehemmten Blutung in den Bauchraum. \gEs{Der Patient
verfällt von einem Moment auf den anderen} und wird
schockig.

\paragraph{Primäre Milzruptur} Die Milz und die Kapsel
reißen gleichzeitig, es kommt zu einem unmittel merkbaren
Schockzustand.

\paragraph{Sekundäre Milzruptur} Die Kapsel reißt erst
einige Zeit nach dem Riß des Milzgewebes \gE{(Latenzzeit)}.
Es kommt zu einem \gEs{zeitlich versetztem} und oft
\gEs{unerwartetem Schockgeschehen}. Mitunter reißt die
Kapsel erst Stunden nach dem eigentlichen Unfall. \gEs{Daher
ist es wichtig jeden Patienten, bei dem auch nur der geringste
begründete Verdacht auf eine Milzverletzung besteht, einer
weiterführenden Untersuchung im Spital zuzuführen!}
\gZ{Mittels Ultraschall kann auf eine Milzblutung untersucht
werden.}

\begin{itemize}
    \item prim\"ar: sofort, Kapsel durchgerissen
    \item sekund\"ar: Sickerblutung, Kapsel anfangs noch intakt
    ${\rightarrow}$ Latenz bis Ri{\ss} der Kapsel
\end{itemize}

\gCave{Cave! {\quotedblbase}\index{Zweizeitige
Milzruptur}Zweizeitige Milzruptur{\textquotedblleft}}



\subsubsection{Weitere interne Verletzungen}

Der Bauchraum beherbergt viele, teilweise sehr gut
durchblutete Organe die bei einer Verletzung zu erheblichen
Blutverlusten führen können. Präklinisch äußern sich diese
Verletzungen durch den durch den Blutverlust hervorgerufenen
Schockzustand. Weitere hinweiszeichen sind Prellmarken oder
Schmerzen in der entsprechenden Gegend.

Grundsätzlich kann man ursächlich gegen diese Verletzung
wenig tun, im Vordergrund steht die Erkennung,
Schockbehandlung und der zügige Transport in ein
Traumazentrum.

\paragraph{Aortenruptur}\index{Aortenruptur}\label{topic:aortenruptur}Durch
einwirkende starke Kräfte kann es zu einem Riß in der Aorta
kommen. Besonders bei Stürzen aus großen Höhen kann dies
vorkommen. Da es sich bei der Aorta um das größte Gefäß des
Körpers handelt ist die Prognose äußerst schlecht.

\paragraph{Nierenruptur}\index{Nierenruptur}\label{topic:nierenruptur}
Die Niere als naturgemäß gut durchblutetes Organ kann im
Falle einer Verletzung auch massive Probleme machen.

\paragraph{Leberuptur} Die Leber ist auch gut durchblutet
und kann im Falle einer Verletzung große Probleme machen.


\subsection{Beckentrauma}

\begin{itemize}
    \item erhebliche Gewalteinwirkung  (Versch\"uttung, \"Uberrolltrauma)
    \item massive  Blutungen (bis 5 l in wenigen  Minuten
    \item Gefahr des Ausblutens (ev. ohne sichtbare  Blutungsquelle!)
    \item Mitverletzung von Organen im kleinen  Becken:
    \begin{itemize}
        \item Dickdarm ${\rightarrow}$ Blutung aus Anus
        \item Harnsystem ${\rightarrow}$ Blut im Harn / Harn in der
        Bauchh\"ohle
        \item (Ausnahme: Spontanruptur der gef\"ullten  Harnblase bei stumpfem
        Bauchtrauma)
        \item Hoden ${\rightarrow}$ H\"amatome am Skrotum und am  Damm
    \end{itemize}
\end{itemize}


\begin{gMassnahmenEnv}
    \paragraph{Spezielle Maßnahmen bei Beckentraumen}
    \begin{itemize}
        \item \gTxMassVitBes
        \begin{itemize}
            \item Lagerung: flach
            \item $O_2$: 10-15\,l
        \end{itemize}
        \item Schockbek\"ampfung
        \item Bei Beckenfraktur:
        \index{Beckengurt}\gEs{Beckengurt}, zur Not mit
        G\"urtel
    \end{itemize}
\end{gMassnahmenEnv}



%%%%%%%%%%%%%%%%%%%%%%%%%%%%%%%%%%%%%%%%%%%%%%%%%%%%%%%%%%%
%%%%%%%%%%%%%%%%%%%%%%%%%%%%%%%%%%%%%%%%%%%%%%%%%%%%%%%%%%%
%%%%%%%%%%%%%%%%%%%%%%%%%%%%%%%%%%%%%%%%%%%%%%%%%%%%%%%%%%%







%%%%%%%%%%%%%%%%%%%%%%%%%%%%%%%%%%%%%%%%%%%%%%%%%%%%%%%%%%%
%%%%%%%%%%%%%%%%%%%%%%%%%%%%%%%%%%%%%%%%%%%%%%%%%%%%%%%%%%%
%%%%%%%%%%%%%%%%%%%%%%%%%%%%%%%%%%%%%%%%%%%%%%%%%%%%%%%%%%%

\section{Polytrauma}

\gAuthNotes{Neufassung}

\marginlabel{Polytrauma}\gDefinition{Gleichzeitige
Verletzung \gEs{mehrerer K\"orperregionen} oder
Organsystemen, von denen \gEs{wenigstens eine} allein
\gEs{oder die Kombination mehrerer} \gEs{lebensgef\"ahrlich}
ist.}

Zum Beispiel:
\begin{itemize}
    \item Bauchtrauma + Oberschenkel; SHT + WS + Thoraxtrauma etc.
    \item gebrochener Finger + gebrochene Zehe ${\rightarrow}$ kein
    Polytrauma
\end{itemize}

\gFazit{Patient gilt als polytraumatisiert  bis zum Beweis 
des Gegenteils!} 

\begin{gMassnahmenEnv}
    \paragraph{Spezielle Maßnahmen bei Polytrauma}
    \begin{itemize}
        \item Schneller Traumacheck gleichzeitig  mit  Sichern der
        Vitalparameter!
        \item Bewu{\ss}tsein
        \item Atmung
        \item Kreislauf: Carotispuls, Pulse peripher  (orientierend!),
        Kapillardurchblutung
        \item gr\"undlicher Traumacheck erst sp\"ater
        \item stabilisierende Ma{\ss}nahmen
        \item Atemwege sichern (Intubation),  ven\"oser  Zugang (Notarzt!)
        \item wenn einigerma{\ss}en stabil, gr\"undlicher  Traumacheck,
        Verletzungsmuster (={\textgreater} zu  erwartende Komplikationen)
    \end{itemize}
\end{gMassnahmenEnv}


\subsubsection{Polytrauma beim Kind}

Der Schockindex ist beim Kind unbrauchbar, da die
Normalwerte andere sind!

Kinder bleiben lange stabil, verfallen dann aber schlagartig!

\subsubsection{Der Unfallort}

\begin{itemize}
    \item Absichern  der Unfallstelle,  m\"oglichst auff\"allig machen (Blaulicht,
 Warnblinkanlage, Scheinwerfer)
 \item bei mehreren Verletzten: Triage
 \item Nachalarmierung : andere  Rettungsmittel, Feuerwehr, Polizei
 \item Eingeklemmte im Fahrzeug stabilisieren,  nicht eigenm\"achtig bergen
 \item Crash -Rettung nur  bei  Fremdgef\"ahrdung, HWS stabilisieren!
\end{itemize}



%%%%%%%%%%%%%%%%%%%%%%%%%%%%%%%%%%%%%%%%%%%%%%%%%%%%%%%%%%%
%%%%%%%%%%%%%%%%%%%%%%%%%%%%%%%%%%%%%%%%%%%%%%%%%%%%%%%%%%%
%%%%%%%%%%%%%%%%%%%%%%%%%%%%%%%%%%%%%%%%%%%%%%%%%%%%%%%%%%%

\section{Extremit\"atenverletzungen}
Gelenke betreffend

Verstauchung (Distorsion )

Verrenkung (Luxation )

Knochenbruch (Fractur )

\subsection{Verstauchung}
\gDefinition{Zerrung / \"Uberdehnung des
Kapsel-/Bandapparates}

\paragraph{Symptome}
\begin{itemize}
    \item Schmerzen
    \item Schwellung
    \item Blutergu{\ss} (H\"amatom )
\end{itemize}



\begin{gMassnahmenEnv}
    \paragraph{Spezielle Massnahmen bei einer Verstauchung}
    \begin{itemize}
        \item ruhigstellen
        \item hochlagern
        \item k\"uhlen
        \item ad Unfallabteilung (mu{\ss} nicht sofort sein)
    \end{itemize}

    
    
\end{gMassnahmenEnv}

\subsection{Verrenkung}

\gDefinition{Teilweise oder vollst\"andige Verrenkung  eines
Gelenks}

\paragraph{Symptome}
\begin{itemize}
    \item Schmerzen
    \item Schwellung
    \item Blutergu{\ss}
    \item abnorme Gelenkstellung (federnd fixiert), bewegungsunf\"ahig
\end{itemize}



\begin{gMassnahmenEnv}
    \paragraph{Spezielle Maßnahmen bei einer Verrenkung}
    
    Stellung belassen
    
    ruhigstellen
    
    KEINE EINRENKVERSUCHE!!!
\end{gMassnahmenEnv}



%%%%%%%%%%%%%%%%%%%%%%%%%%%%%%%%%%%%%%%%%%%%%%%%%%%%%%%%%%%
%%%%%%%%%%%%%%%%%%%%%%%%%%%%%%%%%%%%%%%%%%%%%%%%%%%%%%%%%%%
%%%%%%%%%%%%%%%%%%%%%%%%%%%%%%%%%%%%%%%%%%%%%%%%%%%%%%%%%%%

\subsection{Knochenbr\"uche}

\paragraph{Arten:}

\begin{itemize}
    \item geschlossene Fraktur
    \item offene Fraktur
\end{itemize}

jeweils

\begin{itemize}
    \item komplett
    \item inkomplett
\end{itemize}

\paragraph{Gefahren}

\begin{itemize}
    \item starke Blutungen
    \item ev. bis Schock
    \item begleitende Nervenverletzungen
    \item Fettembolie (Fraktur langer R\"ohrenknochen)
    \item Infektionsgefahr (Osteomyelitis bei offenen  Frakturen)
\end{itemize}

Immer gr\"undlichen Traumacheck  durchf\"uhren, um keine 
Begleitverletzungen zu \"ubersehen!!!


\paragraph{Frakturzeichen} \index{Frakturzeichen}

Es gibt sichere und unsichere Frakturzeichen:

\noindent\begin{tabularx}{\linewidth}[H]{XX}
\toprule
\multicolumn{2}{c}{\gEs{Frakturzeichen}}
\\\midrule
\multicolumn{1}{c}{\gEs{Sichere}}
 &
\multicolumn{1}{c}{\gEs{Unsichere}}
\\\toprule
\begin{itemize}
    \item Abnorme Beweglichkeit
    \emph{("`Pseudogelenk"')}\footnote{"`Pseudogelenk"':
    \emph{"`Wie ein Gelenk"'.} Beim Pesudogelenk hat es den
    Anschein als würde en Gelenk bestehen (das nicht
    dorthon gehört).)}
    \item Fehlstellung
    \item Knochenreiben (Krepitus) \footnote{\so{NICHT}
    absichtlich ausprobieren!}
    \item Sichtbare freie Knochenteile (= offener Bruch)
\end{itemize}
 &
 \begin{itemize}
     \item Schmerz
     \item Schwellung
     \item H\"amatom
     \item gest\"orte Funktion / Bewegungsunf\"ahigkeit
 \end{itemize}
\\\bottomrule
\end{tabularx}
\captionof{table}{Frakturzeichen}

\begin{figure}[h]
    \begin{center}
\includegraphics[width=6cm]{./images/fraktur_fehlstellung_nativ.jpg}
\includegraphics[width=6cm]{./images/fraktur_fehlstellung_xray.jpg}
\caption{Fehlstellung -- In natura und in
Röntgendarstellung. \gSourcePicture{Hauer}\gAuthNotesPicLegalOk{}{}}
\end{center}
\end{figure}

\paragraph{Spezielle Maßnahmen}

Blutstillung

Wundversorgung

ev. NAW-Nachalarmierung (starke Schmerzen, gr\"o{\ss}erer Blutverlust,
Schock, relevanten Begleit\-verletzungen)

Ruhigstellung / Schienung (unter Zug)


%%%%%%%%%%%%%%%%%%%%%%%%%%%%%%%%%%%%%%%%%%%%%%%%%%%%%%%%%%%
%%%%%%%%%%%%%%%%%%%%%%%%%%%%%%%%%%%%%%%%%%%%%%%%%%%%%%%%%%%
%%%%%%%%%%%%%%%%%%%%%%%%%%%%%%%%%%%%%%%%%%%%%%%%%%%%%%%%%%%
%%%%%%%%%%%%%%%%%%%%%%%%%%%%%%%%%%%%%%%%%%%%%%%%%%%%%%%%%%%

\section{Verbrennungen}

\gDefinition{\gAuthNotes{Definition fehlt}}

Die Gr\"o{\ss}e des gesch\"adigten Areals wird in Prozent der
ver\-brannten K\"orperoberfl\"ache angegeben.

\paragraph{Neunerregel}\label{topic:neunerregel}\index{Neunerregel}

Handfl\"ache (des Patienten) = 1~\% der K\"orperoberfl\"ache (KOF)

\gMarried{
\begin{center}
%\includegraphics[width=6.591cm,height=7.03cm]
\gAuthNotes{./images/AASdev20090228-img94.png}

\end{center}
}{
\begin{center}
\begin{tabulary}{\linewidth}[H]{LCC}
\toprule
 & \% &
\\\midrule
Kopf:
 & 9  &
\\
Arm, je:  &  9  &
\\
Thorax vorne: & 18  &
\\
Thorax hinten: & 18 &
\\
Bein, je: & 18 &
\\
Genital:  & 1 &
\\\bottomrule
\end{tabulary}
\end{center}
}

\paragraph{Gradeinteilung}

In Abhängigkeit von der Tiefe der Vebrennung wird das
Ausma{\ss} der Sch\"adigung durch Verbrennung wird in 3 Grade eingeteilt:

\begin{center}
\begin{tabulary}{\linewidth}[H]{CLC}
\toprule
 Grad
 &
 &
\\\midrule
1\,\textdegree
 &
R\"otung, Schmerz, Schwellung
 &
\\
2\,\textdegree
 &
Blasenbildung, Schmerz
 &
\\
3\,\textdegree
 &
Verkohlung, \gT{Verschorfung}, Schmerzen
abhängig von der teife der Schädigung
(Zerstörung von Nerven), Schrumpf\-ung des Gewebes durch
Hitze\-ein\-wirkung & \\\bottomrule
\end{tabulary}
\end{center}

\gCave{Ab \gEs{5\%} verbrannter (mind. 2. Grad) KOF bei
Kindern oder \gEs{10\%} bei Erwachsenen besteht Schockgefahr
(Fl\"ussigkeitsverlust)}


\begin{gMassnahmenEnv}
    \paragraph{Spezielle Ma{\ss}nahmen bei leichter
    Vebrennung}
    
    \begin{itemize}
        \item Eigenschutz!
        \item Kleiderbrand l\"oschen
        \item Kleidung rasch entfernen (sofern nicht  eingeschmolzen!)
        \item mind. 10 min. mit m\"a{\ss}ig kaltem Wasser d
       ie betroffenen Körperregionen spülen  --
       \gEs{NICHT unterk\"uhlen!}
        \item keimfreier, nicht verklebender Verband, locker anlegen
        \item wenn vorhanden: Water-Jel
        \footnote{Streitfrage: Water-Jel: Umstritten}
        \item Schockbek\"ampfung
        \item NAW wenn n\"otig
    \end{itemize}
    \gCave{CAVE Unterkühlung bei zu langer oder zu kalter
    Spülung!}
\end{gMassnahmenEnv}

\begin{gMassnahmenEnv}
    \paragraph{Spezielle Ma{\ss}nahmen bei großflächiger
    Vebrennung}
    \begin{itemize}
        \item \gTxMassVitBes
        \begin{itemize}
            \item $O_2$: 10\gLMin
        \end{itemize}
        \item Kleidung rasch entfernen (sofern nicht  eingeschmolzen!)
        \item mind. 10 min. mit m\"a{\ss}ig kaltem Wasser d
       ie betroffenen Körperregionen spülen  --
       \gEs{NICHT unterk\"uhlen!}
        \item keimfreier, nicht verklebender Verband, locker anlegen
        \item wenn vorhanden: Water-Jel
        \footnote{Streitfrage: Water-Jel: Umstritten}
        \item Schockbek\"ampfung
    \end{itemize}
    \gCave{CAVE Unterkühlung bei zu langer oder zu kalter
    Spülung!}
\end{gMassnahmenEnv}



\paragraph{Was man besser \gE{nicht} tun sollte \ldots}

\begin{itemize}
    \item NEIN: Salbe, Puder, etc. auf frische  Brandwunde
    \item NEIN: Blasen \"offnen
    \item NEIN: festklebende Kleidung gewaltsam  entfernen
    \item NEIN: aluminiumbeschichteten Verband  verwenden (Gefahr der
    Hitzereflexion bei  nicht ausreichend gek\"uhlten  Verbrennungen!)
\end{itemize}



\subsection{Inhalationstrauma}

\paragraph{Leitsymptome}

\begin{itemize}
    \item angebrannte Gesichtshaare
    \item geschwollene, ru{\ss}ige Lippen
    \item Ru{\ss}partikel im Auswurf
    \item Heiserkeit
    \item Stridor
\end{itemize}

${\rightarrow}$ Rauchgasvergiftung m\"oglich!

\begin{gMassnahmenEnv}
    \paragraph{Spezielle Maßnahmen beim Inhalationstrauma}~
    \gAuthNotes{fehlt!}
\end{gMassnahmenEnv}



%%%%%%%%%%%%%%%%%%%%%%%%%%%%%%%%%%%%%%%%%%%%%%%%%%%%%%%%%%%
%%%%%%%%%%%%%%%%%%%%%%%%%%%%%%%%%%%%%%%%%%%%%%%%%%%%%%%%%%%
%%%%%%%%%%%%%%%%%%%%%%%%%%%%%%%%%%%%%%%%%%%%%%%%%%%%%%%%%%%

\section{Erfrierungen}
\gDefinition{...sind Gewebsnekrosen, die durch lokale 
Durchblutungsst\"orungen aufgrund von  K\"alteeinwirkung (Eis, Wind, N\"asse,  K\"alte)
entstanden sind.}

Besonders geff\"ahrdet sind k\"orperferne K\"orperteile wie Zehen,
Finger, Ohren oder Nase.

Einteilung erfolgt in 4 Stadien, diese haben am Notfallort keine
entscheidende Bedeutung f\"ur die zu setzenden Ma{\ss}nahmen. Den
Schweregrad einer Erfrierung kann man erst nach Tagen voll
einsch\"atzen. Deshalb ist es wichtig, die Warnsymptome ernst zu
nehmen!

\paragraph{Stadien}

\gAuthNotes{fehlt!}

\paragraph{Symptome}

\begin{itemize}
    \item zuerst
    \begin{itemize}
        \item Gef\"uhllosigkeit und Bl\"asse
        \item starke Schmerzen und bl\"auliche Verf\"arbung
    \end{itemize}
    \item einige Zeit sp\"ater
    \begin{itemize}
        \item  marmorierte Verf\"arbung der Haut mit  Blasenbildung
        \item Taubheitsgef\"uhl, keine Empfindung bei  Ber\"uhrung
        \item Bewegungsunf\"ahigkeit des betroffenen  K\"orperteils
        \item starke Schmerzen
    \end{itemize}
\end{itemize}


\begin{gMassnahmenEnv}
    \paragraph{Spezielle Ma{\ss}nahmen bei Erfrierungen}~
\begin{itemize}
\item beengende Kleidungsst\"ucke lockern (Einschn\"urung erh\"oht
Risiko einer Erfrierung)
\item Keimfreier, gepolsterter Verband
\item Warme, gezuckerte Getr\"anke -- KEIN Alkohol
\item K\"orperstamm w\"armen, nie das erfrorene K\"orperteil abreiben!!
\end{itemize}
\end{gMassnahmenEnv}




%%%%%%%%%%%%%%%%%%%%%%%%%%%%%%%%%%%%%%%%%%%%%%%%%%%%%%%%%%%
%%%%%%%%%%%%%%%%%%%%%%%%%%%%%%%%%%%%%%%%%%%%%%%%%%%%%%%%%%%
%%%%%%%%%%%%%%%%%%%%%%%%%%%%%%%%%%%%%%%%%%%%%%%%%%%%%%%%%%%

\section{Unf\"alle durch Strom- und Blitzschlag}

\marginpar{\includegraphics[width=\linewidth]{./images/electricity_japan.jpg}
\small\emph{Foto: Fina}}

\gT{Niederspannung}: bis 1000\,Volt (Steckdose:  230 Volt)

\gT{Hochspannung}: \"uber 1000\,Volt, Stromst\"arke
{\textgreater}  3 Ampere (\"Uberlandleitungen) 

Die Verletzungsmuster variieren je nach  Spannung
(Volt) und Stromst\"arke (Ampere).

\begin{itemize}
    \item Reizeffekte (Muskel-/Nervenerregung):
    \begin{itemize}
        \item Verkrampfung ({\quotedblbase}kann nicht
        loslassen{\textquotedblleft})
        \item Sensibilit\"atsst\"orungen
        \item ZNS ${\rightarrow}$ Bewu{\ss}tseinsst\"orung
        \item Herzrhythmusst\"orungen (Kammerflimmern, ev.  als Sp\"atfolge!)
    \end{itemize}
    \item W\"armeentstehung
    \begin{itemize}
        \item Strommarken (Ein-/Austrittsverbrennung)
        \item Lichtbogen: oberfl\"achliche und tiefe  Verbrennungen
    \end{itemize}
\end{itemize}

\begin{itemize}
    \item Gleichstrom
    \item Wechselstrom ${\rightarrow}$ eher  Herzrhythmusst\"orungen
    \item Patient  kann Spannungsquelle ev. nicht  loslassen (Krampf der
    Handmuskulatur)
    \item Auch Defi stellt Strom-Gefahrenquelle  dar!!!
    \item Blitzunfall
    \begin{itemize}
        \item mehrere Millionen (!) Volt
        \item bis 100.000 Ampere
    \end{itemize}
\end{itemize}

\paragraph{Symptome}

\begin{itemize}
    \item Bewu{\ss}tlosigkeit
    \item L\"ahmungen (vor\"ubergehend)
    \item Blitzfiguren auf der Haut
    \item Herzrhythmusst\"orungen m\"oglich
\end{itemize}

\gFazit{Das eigentliche Ausma{\ss} der Verletzung ist oft
{\quotedblbase}unsichtbar{\textquotedblleft}, das der Strom im
K\"orperinneren -- also {\quotedblbase}unterirdisch{\textquotedblleft}
- viel Schaden anrichtet. }

\begin{gMassnahmenEnv}
    \paragraph{Spezielle Ma{\ss}nahmen}
    \begin{itemize}
        \item Eigenschutz!
        \begin{itemize}
            \item (Stromkreis  vorher  unterbrechen, FI bzw. bei Starkstrom durch
            Fachmann)
        \end{itemize}
        \item Genaue  Patientenuntersuchung ${\rightarrow}$ Eintritts{}-  und
        Austrittsmarke
        \item auf ev. Herzrhythmusst\"orungen vorbereitet  sein ${\rightarrow}$
        bis 24h Latenz
        \item Vitalparameter ${\rightarrow}$ entsprechende  symptomatische
        Versorgung
        \item NAW ${\rightarrow}$ Hospitalisierung
        \item Stromverbrennungen wie Verbrennungen  versorgen
    \end{itemize}
\end{gMassnahmenEnv}



